\section{Related Work}
\label{sec:related_work}
% refer to recammaster, wroldforge, trajectorycrafter and other recent related approaches. mainly discuss recent related work. their approach, limitations and difference with ours

\subsection{Camera Controls for Video Generation Models}
The introduction of transformer-based diffusion models~\cite{peebles2023scalable,wan21,moviegen,seedance1,cogvideox} has catalyzed an explosion in video generation research and, subsequently, the development of methods that enable precise video control by leveraging the generative priors inherent in these models~\cite{vid_ctrl_fu20243dtrajmaster,vid_ctrl_guo2024sparsectrl,vid_ctrl_xing2024make,vid_ctrl_yin2023dragnuwa}. These additional control signals encompass a variety of modalities, including user interactions like dragging~\cite{drag_geng2025motion,drag_jeong2024dreammotion,drag_jeong2024vmc,drag_zhao2024motiondirector}, explicit camera coordinates~\cite{camera_traj_he2024cameractrl,camera_traj_wang2024motionctrl,camera_traj_xiao2024trajectory,camera_traj_xu2024camco,camera_traj_yang2024direct,camera_traj_you2024nvs,camera_traj_zhang2025recapture}, and human pose estimation~\cite{human_animate_cheng2025wan,human_animate_gan2025humandit,human_animate_qin2023dancing,human_animate_wang2025unianimate}. Some techniques, particularly in video editing contexts, also use a source video directly as a conditional input~\cite{video_edit_cohen2024slicedit,video_edit_ouyang2024i2vedit,video_edit_wang2025videodirector}.

Within the domain of camera-controlled video generation, numerous models have been trained to utilize camera parameters~\cite{camera_traj_he2024cameractrl,gcd}, such as Plücker embeddings, as conditional inputs for guided text-to-video or image-to-video synthesis. Training-free approaches have also been explored~\cite{watson2022novel}. Many methods advocate for the integration of 3D priors, such as 3D meshes~\cite{bai2025recammaster,ex4d,yu2025trajectorycrafter}, and point tracking~\cite{wang2025epic,jeong2025reangle}, to furnish the model with robust spatial cues. These camera-controlled image-to-video techniques often involve converting the source image into a 3D representation, applying camera-coordinate-based transformations to this render, and subsequently utilizing these generated pseudo-anchors to ground the final video synthesis process. Our work builds upon this foundation by enabling more robust and flexible video re-shooting from monocular inputs, specifically addressing the challenges of texture fidelity and temporal consistency when relying on potentially noisy 3D/4D representations.

\subsection{Video-to-Video Generation and Re-shooting}
Camera-controlled video-to-video generation, also known as video re-shooting, has been extensively investigated in several recent works, including \recammaster{}~\cite{bai2025recammaster}, \exfd{}~\cite{ex4d}, and \trajectorycrafter{}~\cite{yu2025trajectorycrafter}. Methods such as \recammaster{} and GCD~\cite{gcd} often rely on synthetic video data generated via simulators like Kubric~\cite{greff2022kubric} to create video training pairs captured from diverse perspectives. A key limitation of these synthetic data-driven approaches, which our work directly addresses, is their inherent generalization gap: they are often challenging to scale effectively due to the labor-intensive nature of creating diverse synthetic scenarios that accurately reflect real-world complexity, and models trained on them typically fail to generalize to real-world samples or novel object interactions outside the training distribution.

To address this scaling problem and the synthetic-to-real gap, other methods employ 4D point tracking and depth estimation models on in-the-wild videos~\cite{ex4d,jeong2025reangle,chen2025reconstruct}. The projections derived from these models serve as anchors for training the video generation model, analogous to an inpainting task for regions that become occluded or disoccluded. Techniques proposed in works like EPiC~\cite{wang2025epic} and \recammaster{} suggest applying data augmentations to these anchors to mitigate the domain gap between the noisy 3D anchor representations and the high-fidelity source videos. Nonetheless, a fundamental constraint of these methods, which our self-supervised approach aims to overcome, is their reliance on the inherent deficiencies of current 3D/4D point tracking and reconstruction models, and the pervasive scarcity of accurately paired real-world multi-view video data. Methods like \trajectorycrafter\cite{yu2025trajectorycrafter} and SyncCamMaster~\cite{bai2024syncammaster} attempt to mitigate this limitation by training on a heterogeneous mixture of anchor videos and static multi-view datasets.

Concurrent work, such as Reangle-A-Video~\cite{jeong2025reangle}, utilizes point-based warping on in-the-wild videos, combined with stochastic-control guidance, to enhance coherence with the input video. \exfd{} posits that the quality of 4D meshes and anchors can be substantially enhanced by employing watertight meshes. Other approaches, like "Reconstruct, Inpaint, Finetune"~\cite{chen2025reconstruct}, advocate for test-time training on the specific input video to improve coherence and reduce over-reliance on the anchor video.

In contrast, our framework introduces a scalable, domain-agnostic methodology for video re-shooting, generating training pairs directly from monocular videos without relying on complex 3D/4D reconstruction or synthetic data. This self-supervised approach enables our diffusion model to implicitly learn robust 3D reasoning and reliably source high-fidelity textures from the original input. We demonstrate that our architectural adaptations and novel training strategies facilitate strong alignment with desired camera parameters and high coherence with the source video, significantly advancing the capabilities of robust video re-shooting.